\begin{env}[4.2.3]
\label{0.4.2.3-fr}
Sean $\mathcal{M}$, $\mathcal{N}$ dos $\mathcal{A}$-Módulos. Para todo
abierto $U\subset Y$, se tiene por definición
\[
  \Gamma\bigl(\psi^{-1}(U),
    \mathcal{H}\!om_{\mathcal{A}}(\mathcal{M},\mathcal{N})\bigr)
  =
  \operatorname{Hom}_{\mathcal{A}|V}
    (\mathcal{M}|V,\mathcal{N}|V),
\]
donde se ha puesto, para abreviar, $V=\psi^{-1}(U)$. La aplicación
$u\to\Psi_*(u)$ es un homomorfismo
\[
  \operatorname{Hom}_{\mathcal{A}|V}(\mathcal{M}|V,\mathcal{N}|V)
  \longrightarrow
  \operatorname{Hom}_{\mathcal{B}|U}
    (\Psi_*(\mathcal{M})|U,\Psi_*(\mathcal{N})|U)
\]
para las estructuras de $\Gamma(U,\mathcal{B})$-módulos. Como estos
homomorfismos son compatibles con las operaciones de restricción, definen
un homomorfismo canónico funtorial de $\mathcal{B}$-Módulos
\[
  \Psi_*\bigl(\mathcal{H}\!om_{\mathcal{A}}(\mathcal{M},\mathcal{N})\bigr)
  \longrightarrow
  \mathcal{H}\!om_{\mathcal{B}}
    (\Psi_*(\mathcal{M}),\Psi_*(\mathcal{N})).
  \tag{4.2.3.1}\label{0.4.2.3.1-fr}
\]
\end{env}
